\section{结论}

\method{}把确定性显著性场转化为随图像变化的视觉伪时序，同时直接从原始RGB图像读取内容。同一张Sobel场为整圈选择低显著性侧，在周长固定的点数预算内重分配采样点，并确定公共锚点。反向法线和双向闭环独立编码每个轮廓；终止圈通过第二次闭环汇集两侧法线。随后，稀疏增广等高线树从低到高汇总带Tag的轮廓输出，用复制处理分裂，用与排列无关的融合处理汇合。最终叶子只在全局读出时按实际读取区域的显著性排序。与自适应排列或空间分区不同，本文保留等值线分量的分裂与汇合演化，并把局部读取与全局拓扑分开。本文完整定义这一机制，但不声称实证性能优势。学习型显著性、更丰富的原图读取、分支摘要和树并行执行属于可选扩展。
